
\documentclass[11pt]{article}
\usepackage[margin=1in]{geometry}
\usepackage{amsmath,amssymb,booktabs,hyperref,graphicx,parskip,enumitem}
\usepackage{xcolor}
\hypersetup{colorlinks=true,linkcolor=blue!50!black,urlcolor=blue!50!black}
\setlist{nosep}

\title{Replication Report: Effect of Single Atom Platinum (Pt) Doping and Facet on La2Ti2O7\\[2pt]\large OSTI 1981773 \textbar{} Coverage 8/10, Agreement 8/10}
\author{Rick Stevens \and Ollie (AI Assistant)\\\small Argonne National Laboratory}
\date{2026-04-27}

\begin{document}\maketitle

\section{Source paper}
\begin{itemize}
\item \textbf{Title:} Effect of Single Atom Platinum (Pt) Doping and Facet on La2Ti2O7
\item \textbf{Authors:} Galiullin, Shibuya, Yang et al.
\item \textbf{Year:} 2023
\item \textbf{Domain:} Materials / Photocatalysis / DFT
\item \textbf{Identifier:} OSTI 1981773
\end{itemize}


\section{Replication objective}
The central claim of the paper concerns materials / photocatalysis / dft. Our objective was to reproduce the headline numerical/qualitative results within the constraints of the AI-driven Replication Project (24--72\,h, commodity GPU/CPU compute, no author contact). A replication plan is archived at \texttt{1981773-Effect-of-Single-Atom-Platinum-Pt/replication\_plan.pdf}.

\section{Methodology}
We followed the standard Replication Project workflow: ingest the source PDF, extract method and data pointers, draft a replication plan, attempt the smallest end-to-end pipeline that produces a comparable number, then scale up where compute and data permit. Tools used in this replication are catalogued in the meta-paper's computational tools inventory (see \texttt{COMPUTATIONAL\_TOOLS\_INVENTORY.md}).

This paper went through the Tier-Lift pipeline; supplementary notes at \texttt{.openclaw/workspace/24h-progress/batch3-materials/report\_1981773\_LTO.md}, \texttt{.openclaw/workspace/24h-progress/tier-lift-v2.5/1981773/REPORT.md}, \texttt{.openclaw/workspace/24h-progress/tier-lift-v2/1981773.md}.

\subsection*{Paper-directory README excerpt}
\small
\par\medskip\textbf{Effect of Single Atom Platinum (Pt) Doping and Facet Dependent on the Electronic Structure and Light Absorption of Lanthanum Titanium Oxide (La2Ti2O7): A Density Functional Theory Study}\par
\par$\bullet$\ **OSTI ID:** 1981773
\par$\bullet$\ **Rank:** \#4
\par$\bullet$\ **Replication Score:** 10/10

\par\medskip\textbf{\# Why This Paper}\par
This DFT study is fully computational, uses well-defined simulation data, and provides explicit methodological details, allowing for complete automation and quantitative validation by AI.

\par\medskip\textbf{\# Replication Plan}\par
AI would set up and run DFT calculations in VASP or similar software, compute surface energies, band structures, optical absorption spectra, and perform Bader charge analysis, comparing results to those reported.

\par\medskip\textbf{\# Status}\par
\par$\bullet$\ [ ] Paper reviewed
\par$\bullet$\ [ ] Data identified
\par$\bullet$\ [ ] Code implemented
\par$\bullet$\ [ ] Results validated
\normalsize

The replication was driven by an OpenClaw agent loop (Claude Opus / GPT-5 class models) issuing shell, Python, and (where applicable) GPU jobs. Compute usage and wall-time for individual papers are summarised in the meta-paper's resource table; not separately recorded for this paper unless noted in the Tier-Lift artefact. Sub-agents were spawned for replication-plan generation and (for papers in batches 1--4) for tier-lift extension runs.

\section{What was reproduced}
\begin{itemize}
\item Pt cation character (Loewdin +2.60 e vs Ti +1.76 e)
\item Non-magnetic ground state (PBE)
\item (001) facet relaxation, 44 atoms in slab
\item Gap reduction on Pt doping (0.78 rel vs paper 0.69)
\item Absolute redshift (1.77 eV PBE vs paper 2.2 eV)
\item NEW v2.5: Full eps(omega) absorption with Yambo PBE/RPA on bulk, (001) clean slab, (001)+Pt slab
\item NEW v2.5: HSE06 hybrid gap correction = +1.55 eV consistently across all 3 systems
\item NEW v2.5: Pt-induced visible-light absorption enhancement 2.87x (paper Fig 5 qualitative claim)
\item NEW v2.5: Pt sub-gap absorption onset at 0.52 eV (Im eps \textgreater{} 0.2)
\item NEW v2.5: Yambo redshift 2.40 eV agrees with paper to 9 percent (vs 17 percent for v1 JDOS)
\end{itemize}

\section{What was missing or incomplete}
\begin{itemize}
\item Full BSE excitonic absorption (paper itself does not report BSE; PBE/RPA suffices for matching paper claims; left for future tier if exciton physics needed)
\item (010) slab analysis (paper compares (001) vs (010); we focused budget on (001) where paper reports specific quantitative numbers)
\item Bader charge analysis (we used Loewdin instead; both give same qualitative trend)
\end{itemize}
The dominant blocker categories for this paper, mapped onto the meta-paper's 9-category friction taxonomy (\S4), are listed in \S\ref{sec:friction} below.

\section{Quantitative agreement}
Per-quantity numerical comparisons are not separately tabulated in the structured evaluation record for this paper beyond the agreement rationale below; where the rationale references specific numbers, they are reproduced verbatim. A full numerical comparison table is recorded in the per-replication artefacts (see \S\ref{sec:repro}). For papers below 8/8, the agreement rationale identifies the specific quantities that diverge.

\paragraph{Agreement rationale.} Yambo PBE/RPA absorption red-shift slab clean -\textgreater{} Pt = 2.40 eV vs paper 2.2 eV (9 percent diff, vs v1 JDOS 17 percent diff). HSE06 opens the gap by 1.567 eV (bulk), 1.576 eV (clean slab), 1.538 eV (Pt slab) - remarkably consistent \textasciitilde{}1.55 eV oxide gap correction, validating PBE underestimate. HSE06 absolute gaps: bulk 4.491 eV, clean slab 3.833 eV, Pt slab 2.024 eV (paper PBE values 3.8 / 3.2 / 1.0 eV). Visible-light (1.7-3.1 eV) absorption average rises from 0.103 (clean) to 0.296 (Pt) - 2.87x enhancement, directly confirming paper's central qualitative claim about Pt activating LTO for visible-light photocatalysis. Pt mid-gap absorption tail starts at 0.52 eV (Im eps \textgreater{} 0.2 threshold), exactly as expected from the PDOS Pt-d states found in v1.

\section{Score and friction-point tagging}\label{sec:friction}
\begin{itemize}
\item \textbf{Coverage:} 8/10 --- v1 retained: PBE relaxation of bulk + (001) clean and Pt-doped 44-atom slabs, JDOS-based IPA absorption proxy, Loewdin charges (Pt 2.60 vs Ti 1.76), non-magnetic ground state confirmed. NEW IN v2.5: (a) Yambo 5.3 PBE/RPA full eps(omega) (401 omega-points 0-8 eV, dipole-weighted, NC ONCV pseudos, q-\textgreater{}Gamma, 2x2x2 (bulk) / 2x2x1 (slabs) k-mesh) on all three systems; (b) HSE06 hybrid-functional gap spot-check (NC pseudos, q-mesh 1x1x1, ecutfock 80 Ry) on all three systems. Both were the precise gaps that left v2 at 7/7 (epsilon.x had rejected USPP pseudos and HSE06 was deferred).
\item \textbf{Agreement:} 8/10 --- Yambo PBE/RPA absorption red-shift slab clean -\textgreater{} Pt = 2.40 eV vs paper 2.2 eV (9 percent diff, vs v1 JDOS 17 percent diff). HSE06 opens the gap by 1.567 eV (bulk), 1.576 eV (clean slab), 1.538 eV (Pt slab) - remarkably consistent \textasciitilde{}1.55 eV oxide gap correction, validating PBE underestimate. HSE06 absolute gaps: bulk 4.491 eV, clean slab 3.833 eV, Pt slab 2.024 eV (paper PBE values 3.8 / 3.2 / 1.0 eV). Visible-light (1.7-3.1 eV) absorption average rises from 0.103 (clean) to 0.296 (Pt) - 2.87x enhancement, directly confirming paper's central qualitative claim about Pt activating LTO for visible-light photocatalysis. Pt mid-gap absorption tail starts at 0.52 eV (Im eps \textgreater{} 0.2 threshold), exactly as expected from the PDOS Pt-d states found in v1.
\end{itemize}

\noindent\textbf{Friction points encountered (taxonomy tags):}
\begin{itemize}
\item None recorded.
\end{itemize}

\section{Follow-on questions}
\begin{itemize}
\item Not recorded.
\end{itemize}

\section{Reproducibility pointers}\label{sec:repro}
\begin{itemize}
\item \textbf{Paper directory:} \texttt{1981773-Effect-of-Single-Atom-Platinum-Pt}
\item \textbf{Replication artefacts:}
\begin{itemize}
\item \texttt{/Users/stevens/Dropbox/REPLICATE-PROJECT/1981773-Effect-of-Single-Atom-Platinum-Pt/replication}
\end{itemize}
\item \textbf{Replication-dir hint from JSONL:} \texttt{\textasciitilde{}/Dropbox/REPLICATE-PROJECT/1981773-Effect-of-Single-Atom-Platinum-Pt/ + /data/stevens/scratch/1981773-bse/ (uicgpu)}
\item \textbf{Status:} Not recorded
\end{itemize}

To re-run, change directory into the paper directory above and consult its \texttt{README.md} for the entry-point script. Most replications follow the convention \texttt{replication/run\_*.sh} or \texttt{replication/<subdir>/run.py}.

\end{document}
